\documentclass[final,12pt,a4paper]{article}

\usepackage{notes}

\oddsidemargin 0.0in
\textwidth 6.5in
\topmargin 0in
\textheight 9in

\title{Oceanic Cooling Models}
\author{Derrick Hasterok}

\begin{document}
\maketitle

\section{Introduction}
There have been a large number of cooling models developed for the
cooling of the oceanic lithosphere.  They fall into several classes and
although numerical methods are now being employed with greater frequency
to solve more complicated versions with $P$-$T$-$X$-dependent
properties.  However, only the $P$-$T$-$X$-independent versions have
been solved analytically.  This document explores the various solutions
and briefly discusses their applicability to oceanic cooling.  All
solutions solve the transient heat equation without sources,
\begin{equation}
    \pd{T}{t} = \kappa \pd[2]{T}{z},
\end{equation}
where $z$ is depth, $t$ is time, and $T$ is temperature.  The thermal
diffusivity, $\kappa$, is defined by the mantle density, $\rho_{m}$,
specific heat capacity, $C_{P}$, and thermal conductivity, $k$, by
$\kappa = k / \rho_{m} C_{P}$.  To solve this PDE, the depth can be
non-dimensionalized by $\eta = z / 2 \sqrt{\kappa t}$ and temperature by
$\theta = (T - T_{0}) / (T_{m} - T_{0})$.  Thus rewriting the PDE in the
form of an ODE,
\begin{equation}
    -\eta \pd{\theta}{\eta} = \frac{1}{2} \pd[2]{\theta}{\eta},
\end{equation}

Surface heat flow, $q_{0}$ for each model is computed using Fourier's Law,
\begin{equation}
    q_{0}(t) = k\left( \pd{T}{z} \right)_{z = 0}.
    \label{eq-heatflow}
\end{equation}
The subsidence, $s$, is a little more difficult to compute (since it
involves integration) but arises from balancing isostatic columns.  The
density variations result from the thermal contraction the lithosphere
and the infilling density of water above the cooler column.  Subsidence
can be computed using
\begin{equation}
    s(t) = \frac{\rho_{m}}{\rho_{m} - \rho_{w}} \alpha_{V}
        \int_{0}^{L} [ T(z,0) - T(z,t) ] dz,
    \label{eq-subsidence}
\end{equation}
where $\rho_{w}$ and $\rho_{m}$ are seawater and mantle density,
respectively.  The volumetric expansivity is given by $\alpha_{V}$, and
the initial geotherm is $T(z,0)$.  Temperature is integrated from the
surface down to a depth $L$ which represents the depth of compensation,
which is infinite in some formulations.

Geoid anomalies across fracture zones are an additional observable that
has been used to test cooling models, although the data tend to have
greater scatter than bathymetry or heat flow resulting from density
anomalies deeper in the mantle.  The geoid height anomaly is expressed
as,
\begin{equation}
    \Delta N = -\frac{2 \pi G}{g} \left[
        \frac{(\rho_{m} - \rho_{w}) s^2}{2}
        + \alpha_{V} \rho_{m} \int_{0}^{L} z ( T(z,0) - T(z,t) ) dz
        \right].
    \label{eq-geoid}
\end{equation}
The geoid offset across a fracture zone can be calculated as a function
of the constant age offset, $(\Delta N_{2} - \Delta N_{1}) / (t_{2} -
t_{1})$. 


\section{Half-Space Cooling}
Temperature:
\begin{equation}
    T(z,t) = T_{0} + (T_{m} - T_{0})\,
        \erf\left( \frac{z}{2\sqrt{\kappa t}} \right)
\end{equation}

Heat Flow:
\begin{equation}
    q_{0}(t) = \frac{k (T_{m} - T_{0})}{\sqrt{\pi \kappa t}}
\end{equation}

Subsidence:
\begin{equation}
    s(t) = \frac{2 \rho_{m} \alpha_{V} (T_{m} - T_{0})}
        {(\rho_{m} - \rho_{w})}\sqrt{\frac{\kappa t}{\pi}}
\end{equation}

Geoid Anomaly:
\begin{equation}
    \Delta N = 
        \frac{2 \pi G \rho_{m} \alpha_{V} \kappa (T_{m} - T_{0})}{g}
        \left[ 
            1 + \frac{2 \rho_{m} \alpha_{V} (T_{m} - T_{0})}
            {\pi (\rho_{m} - \rho_{w})}
        \right] t
\end{equation}


\section{Plate Model}

The plate model solution by \cite{McKenzie1967} was conducted for a 2-D
plate, however, for plate velocities over $\sim$10~km/m.y. the
temperature differences from a 1-D solution is small ($<$10-20~K) for
time $>$1~m.y.  We'll first start with the 1-D solution.

\subsection{1-D plate cooling}

\begin{equation}
    T(z,t) = T_{0} - (T_{m} - T_{0})
    \left[\frac{z}{L} + \frac{2}{\pi}\sum_{n = 1}^{\infty}
    \frac{1}{n} \sin\left(\frac{n \pi z}{L}\right) \exp\left(\frac{n^2
    \pi^2 \kappa t}{L^2}\right) \right]
\end{equation}
This solution consists of two parts, a constant gradient term that
represents the equilibrium solution and the Fourier series of a
triangle.  When these two parts are combined, the result is the initial
condition.

Heat flow
\begin{equation}
    q_{0}(t) = \frac{k (T_{m} - T_{0})}{L}\left[ 1 + 2 \sum_{n = 1}^{\infty}
        \exp\left( -\frac{n^2 \pi^2 \kappa t}{L^2}\right) \right]
\end{equation}

Subsidence
\begin{equation}
    s(t) = \frac{\rho_{m} \alpha_{V} (T_{m} - T_{0}) L}{2 (\rho_{m} - \rho_{w})}
        \left[ 1 - \frac{8}{\pi^2} \sum_{n = 0}^{\infty} \frac{1}{(2n + 1)^2}
        \exp\left( -\frac{(2n + 1)^2 \pi^2 \kappa t}{L^2} \right) \right]
\end{equation}


\subsection{2-D plate cooling}
The PDE solved in the 2-D case adds a lateral spatial dimension
perpendicular to the ridge axis in the direction of plate motion,
\begin{equation}
    \pd{T}{t} = \kappa \left( \pd[2]{T}{z} + \pd[2]{T}{x} \right),
\end{equation}
where is the Peclet number $R = v L /(2 \kappa)$, defining a parameter
$\beta_{n} = (R^{2} + n^2 \pi^2)^{1/2} - R$

\begin{equation}
    T(z,t) = T_{0} - (T_{m} - T_{0})
    \left[ \frac{z}{L} + \frac{2}{\pi}\sum_{n = 1}^{\infty}
    \frac{1}{n} \sin \left(\frac{n \pi z}{L}\right)
    \exp\left(-\frac{\beta_{n} x}{L}\right)\right]
\end{equation}

Heat flow
\begin{equation}
    q_{0}(t) = \frac{k (T_{m} - T_{0})}{L}
        \left[ 1 + 2 \sum_{n = 1}^{\infty}
        \exp\left( - \frac{\beta_{n} x}{L} \right)\right]
\end{equation}

Subsidence
\begin{equation}
    s(t) = \frac{\rho_{m} \alpha_{V} (T_{m} - T_{0}) L}{\rho_{m} - \rho_{w}}
        \left[ 1 - \frac{8}{\pi^2} \sum_{n = 0}^{\infty} \frac{1}{(2 n + 1)^2}
        \exp\left( -\frac{\beta_{(2n + 1)} x}{L} \right) \right]
\end{equation}



\section{Stefan's Problem}
\subsection{Cooling of a molten mantle}
Another class of cooling models include those based on Stefan's problem.
Stefan's problem is commonly used in geoscience to model the cooling of
a lava lake and was used by \cite{Parker&Oldenburg1973} to describe
oceanic cooling.

with boundary conditions $T = T_{0}$ at $z = 0$ and $T = T_{m}$ at $z =
z_{l}(t)$ where $z_{l}(t)$ is the solidification boundary where $z_{l} = 0$
at $t = 0$.  The solution is only valid for $0 \leq z \leq z_{l}(t)$.
For $z > z_{l}(t)$, $T = T_{m}$.  The solution to this system is given
by
\begin{equation}
    T(z,t) = T_{0} + (T_{m} - T_{0})
        \frac{\erf\left(\frac{z}{2 \sqrt{\kappa t}}\right)}{\erf(\lambda)},
    \label{eq-stefan}
\end{equation}
where $\lambda$ is a constant.  To solve this equation, the depth is
non-dimensionalized by $z/2 \sqrt{\kappa t}$.  The definition of
$\lambda$ is an analogous, $\lambda = z_{l}/2 \sqrt{\kappa t}$.  Since
$\lambda$ is defined by the lower boundary can look at the heat flow
across the solidifying front.  The heat released by solidification at
the lower boundary must be conducted upward.  Hence the heat flowing
away from this boundary is given by
\begin{equation}
    k \left( \pd{T}{z} \right)_{z = z_{l}} = \rho_{m} H_{f} \od{z_{l}}{t}.
\end{equation}
and by substituting the $\od{z_{l}}{t} = \lambda \sqrt{\kappa} /
\sqrt{t}$ into the equation above the thermal gradient at the boundary
is
\begin{equation}
    \left( \pd{T}{z} \right)_{z = z_{l}} =
        \frac{\rho_{m} \sqrt{\kappa} H_{f} \lambda}{k \sqrt{t}}.
    \label{eq-grad1}
\end{equation}
We can solve Eq. \ref{eq-stefan} for the the thermal gradient at the
bottom boundary
\begin{equation}
    \left( \pd{T}{z} \right)_{z = z_{l}} =
        \frac{(T_{m} - T_{0})}{\sqrt{\pi \kappa t}}
        \frac{\exp(-\lambda^2)}{\erf(\lambda)}.
    \label{eq-grad2}
\end{equation}
Substituting Eq. \ref{eq-grad1} into to Eq. \ref{eq-grad2} we find,
\begin{equation}
    \frac{C_{P} (T_{m} - T_{0})}{H_{f} \sqrt{\pi}} = 
        \lambda \exp(\lambda^2)\, \erf(\lambda).
\end{equation}
The constant $\lambda$ can be determined by a root finding method such
as Newton's method with starting guess, $[C_{P}(T_{m} - T_{0})/2 H_{f}
\sqrt{\pi}]^{1/2}$.

The surface heat loss can be determined by $k \left(\pd{T}{z}\right)_{z
= 0}$,
\begin{equation}
    q_{0}(t) = \frac{k (T_{m} - T_{0})}
        {2 \sqrt{\kappa t}} [\erf (\lambda)]^{-1}.
\end{equation}

The subsidence is 
\begin{equation}
    s(t) = \frac{\rho_{m} \alpha_{V}(T_{m} - T_{0})}{(\rho_{m} - \rho_{w})}
        \int_0^\infty \left[
        1 - \frac{\erf \left( \frac{z}{2\sqrt{\kappa t}} \right)}{\erf(\lambda)}
        \right] dz,
\end{equation}
\begin{equation}
    s(t) = \frac{2 \rho_{m} \alpha_{V}(T_{m} - T_{0})}{(\rho_{m} - \rho_{w})}
        \sqrt{\frac{\kappa t}{\pi}}
        \left[ \frac{1 - \exp(-\lambda^{2})}{\erf(\lambda)} \right]
\end{equation}

While \cite{Parker&Oldenburg1973} used a 2-D formulation of Stefan's
problem, the differences are again only when the lithospheric age is
younger than 0.5~m.y. or so.  A larger issue is the appropriate value
for the heat of fusion.  \cite{Parker&Oldenburg1973} uses a value
similar to the latent heat of crystallizing basalt 420~kJ/kg.  However,
since the asthenosphere is only partially molten, and may not even be
molten at older ages, the heat released from must be significantly
smaller.


\subsection{Constant lithospheric heat loss}
Implicit in the development of the plate cooling model is a heat flow at the
base of the plate initially zero and increasing through time to an asymptotic
condition.  However, it seems unlikely that the heat loss from the mantle into
the base of the plate would physically increase with time without a physical
mechanism.  The {\it c}\hspace{1pt}onstant {\it h}\hspace{1pt}eat-flow {\it
a}\hspace{1pt}ssigned {\it b}\hspace{1pt}elow the {\it l}\hspace{1pt}ithospheric
{\it is}\hspace{1pt}otherm (CHABLIS) model, so named by \cite{Doin&Fleitout1996}
seeks a more reasonable assumption of a constant supply of heat provided through
mantle convection.  This class of models is to Stefan's problem as the plate
model is to the half-space cooling method.  This derivation is corrected from
the original version by \citet{Hasterok2013a}.

The chablis model is a solution to the one-dimensional transient heat equation,
\begin{equation}
    \pd{T}{t} = \kappa \pd[2]{T}{z},
    \label{eq:pde}
\end{equation}
initially isothermal at the mantle temperature, $T_m$, and a fixed surface
temperature boundary condition, $T_{0}$.  The lower boundary is moving and
defined by the isotherm at $T_m$.  Heat flow across this lower boundary is also
fixed, $q_l$.

To solve Equation~\ref{eq:pde}, the temperature and depth are
non-dimensionalized by
\begin{align*}
    \theta = & \frac{T - T_{0}}{T_{m} - T_{0}} \\
    \eta = & \frac{z}{4} \sqrt{\frac{\pi}{\kappa t}},
\end{align*}
respectively.  At the lower boundary, we will define the non-dimensional length
to be
\begin{equation}
    \eta_l = \lambda = \frac{z_l}{4} \sqrt{\frac{\pi}{\kappa t}}.
    \label{eq:lambda}
\end{equation}
Thus rewriting the partial differential equation in the form of an ordinary
differential equation,
\begin{equation*}
    -\eta \pd{\theta}{\eta} = \frac{1}{2} \pd[2]{\theta}{\eta}.
    \label{eq:ode}
\end{equation*}
The solution is
\begin{equation}
    T(z,t) = 
    \begin{cases}
        T_{0} + (T_{m} - T_{0})
        {\displaystyle \frac{\erf (\eta)}{\erf (\lambda)} } &
        \text{if $0 \leq z \leq z_{l}(t)$} \\
        T_{m} & \text{if $z > z_{l}(t)$}
    \end{cases}.
    \label{eq:chablisT}
\end{equation}

At first glance, the solution is similar to the classic Stefan's problem
\citep{Turcotte&Schubert2002}; however, the value of $\lambda$ is different for
the case of a solidifying boundary.  To find $\lambda$ in this case, we need to
determine $( \partial T / \partial z)_{z=z_{l}}$ from Equation~\ref{eq:chablisT}
first.  Recall that $\partial (\erf u) / \partial u = 2 \pi^{-1/2} \exp
(-u^2)$.  Hence the thermal gradient at the boundary is
\begin{equation*}
    \left( \pd{T}{z} \right)_{z=z_{l}} = 
        \frac{(T_{m} - T_{0})}{2 \sqrt{\kappa t}}
        \frac{1}{\exp (\lambda^2)\, \erf (\lambda)}.
\end{equation*}
Substituting the basal heat flow boundary condition into the result above, we
find
\begin{equation*}
    \frac{k (T_{m} - T_{0})}{2 q_{l} \sqrt{\kappa t}} =
        \exp(\lambda^2)\, \erf (\lambda).
\end{equation*}
The lithospheric growth is limited to an asymptotic plate thickness, $L = k
(T_{m} - T_{0})/q_{l}$.  If we specify the $L$ rather than the $q_{l}$, then the
solution to $\lambda$ is found by solving a transcendental equation,
\begin{equation}
    \frac{L}{2 \sqrt{\kappa t}} = \exp(\lambda^2)\, \erf (\lambda).
    \label{eq:lambdasoln}
\end{equation}
The value of $\lambda$ can again be found using Newton's method.  Unlike the
classical Stefan's problem, the value of $\lambda$ changes with time.

The surface heat loss can be determined by Fourier's Law, 
\begin{equation}
    q_{0}(t) = k \left(\pd{T(t)}{z}\right)_{z = 0} =
        \frac{k (T_{m} - T_{0})} {2 \sqrt{\kappa t}} \frac{1}{\erf (\lambda)}.
\end{equation}
We can simplify the result above using Equation~\ref{eq:lambdasoln},
\begin{equation}
    q_{0}(t) = \frac{k (T_{m} - T_{0})} {L} \exp (\lambda^2).
\end{equation}
The solution is only valid for $t > 0$ because there is a singularity as $t \to
0$.

At $t \to \infty$, $q_{0}$ must approach the basal heat flow, $q_{l}$.  To show
this condition is true, 
\begin{align*}
    \lim_{t \to \infty} q_{0} & = \lim_{\lambda \to 0} 
        \frac{k (T_{m} - T_{0})}{L} \exp (\lambda^2), \\
        & = \frac{k (T_{m} - T_{0})}{L} = q_l .
\end{align*}

Subsidence represents the buoyancy resulting from the integrated thermal
state within the plate (i.e., $0 \leq z \leq z_{l}(t)$), written
\begin{equation}
    s(t) = \frac{\rho_{m}}{\rho_{m} - \rho_{w}} \alpha
        \int_{0}^{L} [ T_m - T(z,t) ]  dz .
\end{equation}
where $\alpha$ is the thermal expansivity.  Substituting
Equation~\ref{eq:chablisT} and after some rearranging, we find
\begin{equation*}
    s(t) = \frac{\rho_{m} \alpha (T_{m} - T_{0})}{\rho_{m} - \rho_{w}}
        \int_{0}^{z_{l}} \left[1 - 
        \frac{\erf (\eta)}
        {\erf (\lambda)}
        \right] dz .
\end{equation*}
Next, we change coordinates of the integral by substituting the differential
differential, $dz = 4 \sqrt{\kappa t/\pi}\, d\eta$,
\begin{equation}
    s(t) = \frac{4 \rho_{m} \alpha (T_{m} - T_{0})}{\rho_{m} - \rho_{w}}
        \sqrt{\frac{\kappa t}{\pi}}
        \int_{0}^{\lambda} \left[1 - \frac{\erf( \eta)}
        {\erf(\lambda)}\right] d\eta .
\end{equation}
Knowing $\int \erf(x)\, dx = x\, \erf(x) + \exp(-x^2)/\sqrt{\pi}$,
and after some algebra, we find
\begin{equation*}
    s(t) = \frac{4 \rho_{m} \alpha (T_{m} - T_{0})}{(\rho_{m} - \rho_{w})}
        \frac{\sqrt{\kappa t}}{\pi}
        \left[ \frac{1 - \exp(-\lambda^{2})}{\erf(\lambda)} \right] ,
\end{equation*}
which we can further simplify using Equation~\ref{eq:lambdasoln}.  Therefore, we
find the subsidence for the chablis model is given by
\begin{equation}
    s(t) = \frac{8 \rho_{m} \alpha (T_{m} - T_{0}) \kappa t}
        {\pi (\rho_{m} - \rho_{w}) L}
        \left[ \exp (\lambda^{2}) - 1 \right] ,
\end{equation}

%\section{$P$-$T$-$X$-Independent versus -Dependent Models}

\bibliographystyle{elsarticle-harv}
\bibliography{/Users/dhasterok/texmf/tex/latex/ref}

\end{document}
